\documentclass[../cours_spe_trigo.tex]{subfiles}
\begin{document}
\setcounter{chapter}{12}
\chapter{Vecteurs, coordonnées polaires et produit scalaire}
\section{Coordonnées cartésiennes et polaires}
Tout vecteur non nul $\vec u$ du plan peut être caractérisé par sa longueur $r=\|\vec u\|>0$ et un angle de direction $\theta$ (défini modulo $2\pi$). Ses coordonnées sont
\[\boxed{\vec u=(r\cos\theta,r\sin\theta).}\]
\begin{demonstration}
Le point extrémité du vecteur $\vec u$ appartient à un cercle de rayon $r$ centré en son origine. Une homothétie de rapport $r$ transforme le point du cercle unité $(\cos\theta,\sin\theta)$ en $(r\cos\theta,r\sin\theta)$, d'où le résultat. $\square$
\end{demonstration}
\section{Produit scalaire et angle : pourquoi apparaît le cosinus ?}
Dans un repère orthonormé, on \textbf{définit} le produit scalaire de $\vec u=(u_x,u_y)$ et $\vec v=(v_x,v_y)$ par
\[\vec u\cdot\vec v=u_xv_x+u_yv_y.\]
\begin{propriete}
Si $\vec u$ et $\vec v$ sont non nuls et forment un angle géométrique $\gamma\in[0,\pi]$, alors
\[\boxed{\vec u\cdot\vec v=\|\vec u\|\,\|\vec v\|\cos\gamma.}\]
\end{propriete}
\begin{demonstration}
D'une part, développer la norme au carré :
\[
\|\vec u-\vec v\|^2=\|\vec u\|^2+\|\vec v\|^2-2\vec u\cdot\vec v.
\]
D'autre part, la loi des cosinus appliquée au triangle défini par $\vec u$, $\vec v$ et leur différence donne
\[
\|\vec u-\vec v\|^2=\|\vec u\|^2+\|\vec v\|^2-2\|\vec u\|\|\vec v\|\cos\gamma.
\]
Comparer les deux membres démontre la formule du produit scalaire. $\square$
\end{demonstration}
\textbf{Application.} $\vec u=(1,2)$ et $\vec v=(2,-1)$ vérifient $\vec u\cdot\vec v=2-2=0$. Comme ils sont non nuls, $\cos\gamma=0$, donc $\gamma=\pi/2$ : ils sont orthogonaux.
\end{document}
